\subsection{比例した $b^2$ 結合の標準形と非零横結合の全排除}
\label{subsec:cubic-third-root-nonunit-proportional-open}
\begin{proposition}[比例self結合の二つの必要標準形]
\label{prop:cubic-third-root-nonunit-proportional-self}
非単位重複rootの $b^2$ 結合が両非零比例なら、$p=x,r=\alpha x$、$\alpha\ne0$ とできる。
自己対称rank2 blockの関係は
\[
 xc+\alpha cx=0,\qquad
 b^2+u xc+v yc-\frac\omega\alpha(u cx+v cy)=0.
\]
$v\ne0$ なら $y$ を変更して $u=0,v=1$ とできる。
$v=0$ なら整域性から $u\ne0,\alpha^2+\omega\ne0$ が必要であり、$u=1$ に正規化できる。
\end{proposition}
\begin{proof}
前節と同じrank2対称分解を使う。
$P=(1,0,0,\alpha,0)^t$ とし、$V_{b^2}=1$ とすると、
二つの主blockの比例 $-\omega$ により
$V=(u,v,1,-\omega u/\alpha,-\omega v/\alpha)^t$ となる。
$v\ne0$ なら $y'=ux+vy$ と選べる。
$v=0$ で $u=0$ 又は $\alpha^2=-\omega$ なら、二関係から純 $b^2$ 関係が得られ、整域性に反する。
残る $u$ は対角scalar変更で1にできる。
\end{proof}

\begin{theorem}[非退化selfの非零横結合全排除]
\label{thm:cubic-third-root-nonunit-proportional-open-empty}
上の $u=0,v=1$ 標準形で、$m=a_{10},n=a_{18},T=a_{21}$ とする。
$n\ne0$ の全層はH3を与えない。
$n=0$ の候補も、次の二領域に限られる：
\[
 T=0,\quad\alpha=-\omega m^2,\quad m\ne0;
 \qquad
 T\ne0,\quad\alpha=m=-\omega^2.
\]
この二領域、$v=0$、および両結合零は、本定理では排除しない。
\end{theorem}
\begin{proof}
元middleの最後二列への射影には、$n^2$ と $m^2+\alpha\omega^2$ のminorがある。
$n\ne0$ ならこの射影はrank2であり、$c^2$ 行の非零 $xy,yx$ 成分から $T\ne0$ となる。
$n=T=1$ としてよい。
実際、$\gamma^3=n/T$ とし、$x,y,c$ をそれぞれ
$n/\gamma,1/\gamma,\gamma$ 倍、$b$ を固定し、topを $1/n$ 倍すると、
self標準形と $\alpha,m$ を保ってこの正規化を得る。

$yb,by$ 行の最後二列は定数単元行列なので、rank2条件は八つの $U^3b$ 係数に関する12線型式となる。
最初の8斉次式の行列式は
\[
 \omega(\alpha+1)(\alpha+\omega)(\alpha+\omega^2)^2.
\]
これが非零なら八係数が零となり、残る $c^2$ 行の非零項に矛盾する。
従って $\alpha=-1,-\omega,-\omega^2$ の三値に限る。
各値で $\mathbb Q(\omega)[m]$ 上のパラメータ定数単元だけによる消去を行う。
$\alpha=-1,-\omega$ では8単元pivot、残余多項式はそれぞれ $m-\omega,m-1$ の非零定数倍となる。
その二点を元22tensorから再構成すると、どちらも三次商次元22であり、AS4値20に反する。

$\alpha=-\omega^2$ では7単元pivotで全残余が零、$a_{15}=Z$ が全 $m$ で自由に残る。
$l=m+\omega^2$、$Y=y+lx$ と置くと、全元関係の像は正確に
\begin{gather*}
 xc=\omega^2cx,\qquad b^2+Yc+\omega^2cY-2lxc=0,\\
 bY=\omega^2Yb,\qquad Yb-\omega^2xb-bx-c^2=0,\\
 bc=x^2-\omega xY-ZY^2,\qquad
 cb=\omega^2x^2-Yx-\omega^2ZY^2
\end{gather*}
となる。元最高tensorの画像を表示六関係で掛け戻し、
unit minorと係数体以外の分母がないことも確認する。
従ってこれは一般点だけの同定でなく、全 $l,Z$ の必要表示である。

この二変数表示の次数3,4商を全関係から構成すると、
正常核の稠密開集合の行列は $41\times7$、正常消去の7個のpivotは全て係数体の非零定数となる。
関係pivotと分母の全非定数因子は
\[
 Z,\quad Zl+\omega/2,\quad Z-\omega^2,\quad Z-2\omega^2
\]
だけに支持される。従ってその積が非零なら正常核は零。
四つの曲線を新しく一変数関数体上で計算したとき、
全関係・正常pivotと分母を含む例外多項式の平方因子を除いたものは
\[
\begin{array}{c|c|c}
\text{曲線}&\text{変数}&\text{例外多項式}\\\hline
Z=0&l&1\\
Z=\omega^2&l&l+\omega^2/2\\
Z=2\omega^2&l&l+\omega^2/4\\
l=-\omega/(2Z)&Z&Z(Z-\omega^2)(Z-2\omega^2).
\end{array}
\]
最後の曲線では $Z\ne0$ が座標条件である。
従って例外点は二つだけであり、いずれも三次商次元22となってAS4でない。
曲線上では実際の商の正常核全体を計算しており、固定部分が8次元へ増える場合も捨てていない。
これで全 $n\ne0$ を排除した。

最後に $n=0$ とする。
最後二列への射影がrank2なら、$yb,by$ 行は右二成分が零なので全行が零となる。
これにより
$a_1=a_2=a_5=a_6=a_8=a_{15}=0$、さらに全関係の $y^2$ 成分も零となる。
二次提示を $(x,c)$ で商にすると $k\langle y,b\rangle/(b^2)$ となり、
次数15の1597がAS4の816を超える。
従って射影rankは1でなければならない。
二つのminorから
\[
 m^2+\alpha\omega^2=0,\qquad T(m+\omega^2)=0
\]
が必要となり、記載した二つの残余に限られる。

公開再現器は \texttt{verify\_cubic\_third\_root\_nonunit\_}\\
\texttt{proportional\_open.py}。
パラメータを残した全線型消去、元画像の恒等式、四曲線と二点を標数0で再生する。
\end{proof}
